% header.tex -- zentrale Preamble / Einstellungen
% Lade in Masterarbeit.tex zuerst die Klasse, dann diese Datei:
% \documentclass[listof=totoc,fontsize=12pt,numbers=noenddot]{classTimo}
% % header.tex -- zentrale Preamble / Einstellungen
% Lade in Masterarbeit.tex zuerst die Klasse, dann diese Datei:
% \documentclass[listof=totoc,fontsize=12pt,numbers=noenddot]{classTimo}
% % header.tex -- zentrale Preamble / Einstellungen
% Lade in Masterarbeit.tex zuerst die Klasse, dann diese Datei:
% \documentclass[listof=totoc,fontsize=12pt,numbers=noenddot]{classTimo}
% % header.tex -- zentrale Preamble / Einstellungen
% Lade in Masterarbeit.tex zuerst die Klasse, dann diese Datei:
% \documentclass[listof=totoc,fontsize=12pt,numbers=noenddot]{classTimo}
% \input{header.tex}

% --------------------------
% Projektpakete (extern, bleiben als \usepackage)
% --------------------------
% Verzeichnisse, Titelblatt, util bleiben als separate Dateien/Pakete
\usepackage{Verzeichnisse}
\usepackage{Titelblatt}
\usepackage{util}

% --------------------------
% Literaturverzeichnis (biblatex)
% --------------------------
% https://www.overleaf.com/learn/latex/Biblatex_citation_styles
\usepackage[style=numeric, 
sorting=none,        % none --> nach reihenfolge der zitate im text
backend=bibtex,      % or backend=biber
bibencoding=utf8,    % .bib file encoding
maxbibnames=99       % give full list of authors (up to 99)
]{biblatex}
\setlength{\bibitemsep}{0.5\baselineskip} % Abstand zwischen zwei Quellen im Litver.
\bibliography{BibTeX_database_FSDR} % oder \addbibresource{Referenzen.bib} bei backend=biber

% (Optional: Beispiele für xpatch-Einsätze aus der alten Vorlage)
% \usepackage{xpatch}
% \xpretobibmacro{url}{\begingroup\color{red}}{}{}
% \xapptobibmacro{url}{\endgroup}{}

% --------------------------
% Seitenlayout / Kopf-/Fußzeile Einstellungen
% (Hinweis: viele dieser Grundfunktionen liefert classTimo bereits;
%  hier werden nur die projekt-spezifischen Anpassungen aus der alten Preamble übernommen)
% --------------------------

% Abstand vor/nach chapter-Überschrift
\renewcommand*\chapterheadstartvskip{\vspace{0cm}} % Abstand vor chapter Überschrift
\renewcommand*\chapterheadendvskip{\vspace{0cm}}   % Abstand nach chapter Überschrift

% Kopf-/Fußzeilen-Setup (scrlayer-scrpage verwendet classTimo bereits; wir überschreiben hier projekt-spezifisch)
\clearpairofpagestyles
\automark{chapter} % was soll in headmark stehen
\ihead[]{\headmark}  % i = links, c = mittig, o = rechts / headmark: Überschiften
\cfoot*{\pagemark} % [] Überschiftseiten, {} normale Seiten, * für beide
%\ohead*{\includegraphics[width=0.2\textwidth]{./Bilder/Stihl_header}}
\addtokomafont{pagehead}{\bfseries} % upshape = Kursiv, bfseries = Fett
\addtokomafont{pagenumber}{\bfseries}
% Auskommentierte zusätzliche Kopf-/Fuß-Optionen aus Vorlage (nur bei Bedarf aktivieren):
% %\pagestyle{scrheadings}
% %\clearscrheadfoot
% %\setheadsepline{0.5pt}
% %\setfootsepline{0.5pt}
% %\renewcommand*\chapterpagestyle{scrheadings}

% --------------------------
% (Auskommentierte) Verzeichnis-Anpassungen aus Vorlage
% Diese Einstellungen betreffen tocloft/tocbasic; da sie im Original auskommentiert waren,
% belasse ich sie hier kommentiert. Aktiviere nur falls benötigt.
% --------------------------
% %\renewcommand{\cftfigpresnum}{Abb. }	% Anpassen des Abbildungsverzeichnises
% %\renewcommand{\cfttabpresnum}{Tab. }
% %\renewcommand{\cftfigaftersnum}{:}
% %\renewcommand{\cfttabaftersnum}{:}
% %\setlength{\cftfignumwidth}{2cm}
% %\setlength{\cfttabnumwidth}{2cm}
% %\setlength{\cftfigindent}{0cm}
% %\setlength{\cfttabindent}{0cm}

% --------------------------
% Kurze Makros und Farben (aus alter Preamble)
% --------------------------
\newcommand\LM{\ensuremath{\mathit{LM}}}
\newcommand\IS{\ensuremath{\mathit{IS}}}

\usepackage{color}
\definecolor{FE_blue}{rgb}{0,0.5686,0.8627} % #0091DC
\definecolor{FE_blue2}{rgb}{0.8627,0.9216,0.9647} % #DCEBF6
\definecolor{FE_gray}{rgb}{0.8471,0.8627,0.8824} % 
\definecolor{AleeRed}{rgb}{0.5,0,0}

% --------------------------
% Mathematische Theoreme (aus Vorlage)
% --------------------------
% auskommentierte Varianten belassen; aktive Definition übernimmt shadethm-Umgebung
% %\usepackage{thm-autoref}
% %\newtheorem{satz}{Satz}
\newshadetheorem{satz}{Satz}
\newcommand{\shadesatzname}{Satz}
\newenvironment{thm}[1][]{%
	\definecolor{shadethmcolor}{FE_gray}%
	\definecolor{shaderulecolor}{rgb}{1.0,0.0,0.0}%
	\setlength{\shadeboxrule}{1pt}%
	\begin{sthm}[#1]%
	}{\end{sthm}}

% --------------------------
% hyperref (aus alter Vorlage) - am Ende der Pakete laden
% --------------------------
\usepackage[
pdftitle={Masterarbeit Festo},
pdfauthor={Steffen Sieder},
pdfsubject={Modellbasierte Funktionsentwicklung von},
pdfkeywords={},
pdftex,
bookmarks,
bookmarksopen,
breaklinks,
colorlinks=TRUE,
hyperindex=TRUE,
plainpages=FALSE,
pdfstartview=FitH,
bookmarksopenlevel=1,
bookmarksnumbered=true,
citecolor=black,
urlcolor=black,
linkcolor=black
]{hyperref}

% --------------------------
% autoref Namen (deutsch)
% --------------------------
\addto\extrasngerman{%
	\renewcommand{\sectionautorefname}{Abschnitt}%
	\renewcommand{\subsectionautorefname}{Abschnitt}%
	\renewcommand{\subsubsectionautorefname}{Abschnitt}%
}

% --------------------------
% Gleichungsnummerierung in Klammern (aus Vorlage)
% --------------------------
\makeatletter
\let\oldtheequation\theequation
\renewcommand\tagform@[1]{\maketag@@@{\ignorespaces#1\unskip\@@italiccorr}} %damit bei Benennung der Gleichung keine Doppelklammer entsteht
\renewcommand\theequation{(\oldtheequation)}
\makeatother

% --------------------------
% TikZ / Styles (auskommentiert in Vorlage) - hier nur als Kommentar
% --------------------------
% %\usepackage{tikz}
% %\usetikzlibrary{shapes,arrows,fit,calc,positioning,automata}
% %\tikzstyle{block} = [draw, fill=blue!20, rectangle, minimum height=3em, minimum width=\blockwidth]

% --------------------------
% Listings-Einstellungen (aus alter Preamble)
% --------------------------
\definecolor{dkgreen}{rgb}{0,0.6,0}
\definecolor{gray}{rgb}{0.5,0.5,0.5}
\definecolor{mauve}{rgb}{0.58,0,0.82}
\lstset{frame=none,
	aboveskip=3mm,
	belowskip=5mm,
	captionpos=b,
	showstringspaces=false,
	columns=flexible,
	basicstyle={\small\ttfamily},
	numbers=left,
	numberstyle=\tiny\color{gray},
	numbersep=5pt,
	keywordstyle=\color{blue},
	commentstyle=\color{dkgreen},
	stringstyle=\color{mauve},
	breaklines=true,
	breakatwhitespace=true,
	tabsize=3,
	inputpath = ./Skripte
}
\lstset{literate=%
	{Ö}{{\"O}}1
	{Ä}{{\"A}}1
	{Ü}{{\"U}}1
	{ß}{{\ss}}1
	{ü}{{\"u}}1
	{ä}{{\"a}}1
	{ö}{{\"o}}1
}
\renewcommand{\lstlistingname}{Skript}
\renewcommand{\lstlistlistingname}{Skriptverzeichnis}

% --------------------------
% Aufzählungs- und Absatz-Einstellungen sowie Inhalts-/Nummerierungstiefe
% --------------------------
\renewcommand{\labelenumi}{\arabic{enumi}.}

\setlength{\parindent}{0pt} % Einrücken nach Bildern verhindern
\setlength{\parskip}{12pt}	% globale Länge festlegen
\setcounter{tocdepth}{2}	% Inhaltsverzeichnis-Tiefe
\setcounter{secnumdepth}{2} % Nummerierungstiefe

% --------------------------
% Ende header.tex (vorläufig)
% --------------------------


% --------------------------
% Projektpakete (extern, bleiben als \usepackage)
% --------------------------
% Verzeichnisse, Titelblatt, util bleiben als separate Dateien/Pakete
\usepackage{Verzeichnisse}
\usepackage{Titelblatt}
\usepackage{util}

% --------------------------
% Literaturverzeichnis (biblatex)
% --------------------------
% https://www.overleaf.com/learn/latex/Biblatex_citation_styles
\usepackage[style=numeric, 
sorting=none,        % none --> nach reihenfolge der zitate im text
backend=bibtex,      % or backend=biber
bibencoding=utf8,    % .bib file encoding
maxbibnames=99       % give full list of authors (up to 99)
]{biblatex}
\setlength{\bibitemsep}{0.5\baselineskip} % Abstand zwischen zwei Quellen im Litver.
\bibliography{BibTeX_database_FSDR} % oder \addbibresource{Referenzen.bib} bei backend=biber

% (Optional: Beispiele für xpatch-Einsätze aus der alten Vorlage)
% \usepackage{xpatch}
% \xpretobibmacro{url}{\begingroup\color{red}}{}{}
% \xapptobibmacro{url}{\endgroup}{}

% --------------------------
% Seitenlayout / Kopf-/Fußzeile Einstellungen
% (Hinweis: viele dieser Grundfunktionen liefert classTimo bereits;
%  hier werden nur die projekt-spezifischen Anpassungen aus der alten Preamble übernommen)
% --------------------------

% Abstand vor/nach chapter-Überschrift
\renewcommand*\chapterheadstartvskip{\vspace{0cm}} % Abstand vor chapter Überschrift
\renewcommand*\chapterheadendvskip{\vspace{0cm}}   % Abstand nach chapter Überschrift

% Kopf-/Fußzeilen-Setup (scrlayer-scrpage verwendet classTimo bereits; wir überschreiben hier projekt-spezifisch)
\clearpairofpagestyles
\automark{chapter} % was soll in headmark stehen
\ihead[]{\headmark}  % i = links, c = mittig, o = rechts / headmark: Überschiften
\cfoot*{\pagemark} % [] Überschiftseiten, {} normale Seiten, * für beide
%\ohead*{\includegraphics[width=0.2\textwidth]{./Bilder/Stihl_header}}
\addtokomafont{pagehead}{\bfseries} % upshape = Kursiv, bfseries = Fett
\addtokomafont{pagenumber}{\bfseries}
% Auskommentierte zusätzliche Kopf-/Fuß-Optionen aus Vorlage (nur bei Bedarf aktivieren):
% %\pagestyle{scrheadings}
% %\clearscrheadfoot
% %\setheadsepline{0.5pt}
% %\setfootsepline{0.5pt}
% %\renewcommand*\chapterpagestyle{scrheadings}

% --------------------------
% (Auskommentierte) Verzeichnis-Anpassungen aus Vorlage
% Diese Einstellungen betreffen tocloft/tocbasic; da sie im Original auskommentiert waren,
% belasse ich sie hier kommentiert. Aktiviere nur falls benötigt.
% --------------------------
% %\renewcommand{\cftfigpresnum}{Abb. }	% Anpassen des Abbildungsverzeichnises
% %\renewcommand{\cfttabpresnum}{Tab. }
% %\renewcommand{\cftfigaftersnum}{:}
% %\renewcommand{\cfttabaftersnum}{:}
% %\setlength{\cftfignumwidth}{2cm}
% %\setlength{\cfttabnumwidth}{2cm}
% %\setlength{\cftfigindent}{0cm}
% %\setlength{\cfttabindent}{0cm}

% --------------------------
% Kurze Makros und Farben (aus alter Preamble)
% --------------------------
\newcommand\LM{\ensuremath{\mathit{LM}}}
\newcommand\IS{\ensuremath{\mathit{IS}}}

\usepackage{color}
\definecolor{FE_blue}{rgb}{0,0.5686,0.8627} % #0091DC
\definecolor{FE_blue2}{rgb}{0.8627,0.9216,0.9647} % #DCEBF6
\definecolor{FE_gray}{rgb}{0.8471,0.8627,0.8824} % 
\definecolor{AleeRed}{rgb}{0.5,0,0}

% --------------------------
% Mathematische Theoreme (aus Vorlage)
% --------------------------
% auskommentierte Varianten belassen; aktive Definition übernimmt shadethm-Umgebung
% %\usepackage{thm-autoref}
% %\newtheorem{satz}{Satz}
\newshadetheorem{satz}{Satz}
\newcommand{\shadesatzname}{Satz}
\newenvironment{thm}[1][]{%
	\definecolor{shadethmcolor}{FE_gray}%
	\definecolor{shaderulecolor}{rgb}{1.0,0.0,0.0}%
	\setlength{\shadeboxrule}{1pt}%
	\begin{sthm}[#1]%
	}{\end{sthm}}

% --------------------------
% hyperref (aus alter Vorlage) - am Ende der Pakete laden
% --------------------------
\usepackage[
pdftitle={Masterarbeit Festo},
pdfauthor={Steffen Sieder},
pdfsubject={Modellbasierte Funktionsentwicklung von},
pdfkeywords={},
pdftex,
bookmarks,
bookmarksopen,
breaklinks,
colorlinks=TRUE,
hyperindex=TRUE,
plainpages=FALSE,
pdfstartview=FitH,
bookmarksopenlevel=1,
bookmarksnumbered=true,
citecolor=black,
urlcolor=black,
linkcolor=black
]{hyperref}

% --------------------------
% autoref Namen (deutsch)
% --------------------------
\addto\extrasngerman{%
	\renewcommand{\sectionautorefname}{Abschnitt}%
	\renewcommand{\subsectionautorefname}{Abschnitt}%
	\renewcommand{\subsubsectionautorefname}{Abschnitt}%
}

% --------------------------
% Gleichungsnummerierung in Klammern (aus Vorlage)
% --------------------------
\makeatletter
\let\oldtheequation\theequation
\renewcommand\tagform@[1]{\maketag@@@{\ignorespaces#1\unskip\@@italiccorr}} %damit bei Benennung der Gleichung keine Doppelklammer entsteht
\renewcommand\theequation{(\oldtheequation)}
\makeatother

% --------------------------
% TikZ / Styles (auskommentiert in Vorlage) - hier nur als Kommentar
% --------------------------
% %\usepackage{tikz}
% %\usetikzlibrary{shapes,arrows,fit,calc,positioning,automata}
% %\tikzstyle{block} = [draw, fill=blue!20, rectangle, minimum height=3em, minimum width=\blockwidth]

% --------------------------
% Listings-Einstellungen (aus alter Preamble)
% --------------------------
\definecolor{dkgreen}{rgb}{0,0.6,0}
\definecolor{gray}{rgb}{0.5,0.5,0.5}
\definecolor{mauve}{rgb}{0.58,0,0.82}
\lstset{frame=none,
	aboveskip=3mm,
	belowskip=5mm,
	captionpos=b,
	showstringspaces=false,
	columns=flexible,
	basicstyle={\small\ttfamily},
	numbers=left,
	numberstyle=\tiny\color{gray},
	numbersep=5pt,
	keywordstyle=\color{blue},
	commentstyle=\color{dkgreen},
	stringstyle=\color{mauve},
	breaklines=true,
	breakatwhitespace=true,
	tabsize=3,
	inputpath = ./Skripte
}
\lstset{literate=%
	{Ö}{{\"O}}1
	{Ä}{{\"A}}1
	{Ü}{{\"U}}1
	{ß}{{\ss}}1
	{ü}{{\"u}}1
	{ä}{{\"a}}1
	{ö}{{\"o}}1
}
\renewcommand{\lstlistingname}{Skript}
\renewcommand{\lstlistlistingname}{Skriptverzeichnis}

% --------------------------
% Aufzählungs- und Absatz-Einstellungen sowie Inhalts-/Nummerierungstiefe
% --------------------------
\renewcommand{\labelenumi}{\arabic{enumi}.}

\setlength{\parindent}{0pt} % Einrücken nach Bildern verhindern
\setlength{\parskip}{12pt}	% globale Länge festlegen
\setcounter{tocdepth}{2}	% Inhaltsverzeichnis-Tiefe
\setcounter{secnumdepth}{2} % Nummerierungstiefe

% --------------------------
% Ende header.tex (vorläufig)
% --------------------------


% --------------------------
% Projektpakete (extern, bleiben als \usepackage)
% --------------------------
% Verzeichnisse, Titelblatt, util bleiben als separate Dateien/Pakete
\usepackage{Verzeichnisse}
\usepackage{Titelblatt}
\usepackage{util}

% --------------------------
% Literaturverzeichnis (biblatex)
% --------------------------
% https://www.overleaf.com/learn/latex/Biblatex_citation_styles
\usepackage[style=numeric, 
sorting=none,        % none --> nach reihenfolge der zitate im text
backend=bibtex,      % or backend=biber
bibencoding=utf8,    % .bib file encoding
maxbibnames=99       % give full list of authors (up to 99)
]{biblatex}
\setlength{\bibitemsep}{0.5\baselineskip} % Abstand zwischen zwei Quellen im Litver.
\bibliography{BibTeX_database_FSDR} % oder \addbibresource{Referenzen.bib} bei backend=biber

% (Optional: Beispiele für xpatch-Einsätze aus der alten Vorlage)
% \usepackage{xpatch}
% \xpretobibmacro{url}{\begingroup\color{red}}{}{}
% \xapptobibmacro{url}{\endgroup}{}

% --------------------------
% Seitenlayout / Kopf-/Fußzeile Einstellungen
% (Hinweis: viele dieser Grundfunktionen liefert classTimo bereits;
%  hier werden nur die projekt-spezifischen Anpassungen aus der alten Preamble übernommen)
% --------------------------

% Abstand vor/nach chapter-Überschrift
\renewcommand*\chapterheadstartvskip{\vspace{0cm}} % Abstand vor chapter Überschrift
\renewcommand*\chapterheadendvskip{\vspace{0cm}}   % Abstand nach chapter Überschrift

% Kopf-/Fußzeilen-Setup (scrlayer-scrpage verwendet classTimo bereits; wir überschreiben hier projekt-spezifisch)
\clearpairofpagestyles
\automark{chapter} % was soll in headmark stehen
\ihead[]{\headmark}  % i = links, c = mittig, o = rechts / headmark: Überschiften
\cfoot*{\pagemark} % [] Überschiftseiten, {} normale Seiten, * für beide
%\ohead*{\includegraphics[width=0.2\textwidth]{./Bilder/Stihl_header}}
\addtokomafont{pagehead}{\bfseries} % upshape = Kursiv, bfseries = Fett
\addtokomafont{pagenumber}{\bfseries}
% Auskommentierte zusätzliche Kopf-/Fuß-Optionen aus Vorlage (nur bei Bedarf aktivieren):
% %\pagestyle{scrheadings}
% %\clearscrheadfoot
% %\setheadsepline{0.5pt}
% %\setfootsepline{0.5pt}
% %\renewcommand*\chapterpagestyle{scrheadings}

% --------------------------
% (Auskommentierte) Verzeichnis-Anpassungen aus Vorlage
% Diese Einstellungen betreffen tocloft/tocbasic; da sie im Original auskommentiert waren,
% belasse ich sie hier kommentiert. Aktiviere nur falls benötigt.
% --------------------------
% %\renewcommand{\cftfigpresnum}{Abb. }	% Anpassen des Abbildungsverzeichnises
% %\renewcommand{\cfttabpresnum}{Tab. }
% %\renewcommand{\cftfigaftersnum}{:}
% %\renewcommand{\cfttabaftersnum}{:}
% %\setlength{\cftfignumwidth}{2cm}
% %\setlength{\cfttabnumwidth}{2cm}
% %\setlength{\cftfigindent}{0cm}
% %\setlength{\cfttabindent}{0cm}

% --------------------------
% Kurze Makros und Farben (aus alter Preamble)
% --------------------------
\newcommand\LM{\ensuremath{\mathit{LM}}}
\newcommand\IS{\ensuremath{\mathit{IS}}}

\usepackage{color}
\definecolor{FE_blue}{rgb}{0,0.5686,0.8627} % #0091DC
\definecolor{FE_blue2}{rgb}{0.8627,0.9216,0.9647} % #DCEBF6
\definecolor{FE_gray}{rgb}{0.8471,0.8627,0.8824} % 
\definecolor{AleeRed}{rgb}{0.5,0,0}

% --------------------------
% Mathematische Theoreme (aus Vorlage)
% --------------------------
% auskommentierte Varianten belassen; aktive Definition übernimmt shadethm-Umgebung
% %\usepackage{thm-autoref}
% %\newtheorem{satz}{Satz}
\newshadetheorem{satz}{Satz}
\newcommand{\shadesatzname}{Satz}
\newenvironment{thm}[1][]{%
	\definecolor{shadethmcolor}{FE_gray}%
	\definecolor{shaderulecolor}{rgb}{1.0,0.0,0.0}%
	\setlength{\shadeboxrule}{1pt}%
	\begin{sthm}[#1]%
	}{\end{sthm}}

% --------------------------
% hyperref (aus alter Vorlage) - am Ende der Pakete laden
% --------------------------
\usepackage[
pdftitle={Masterarbeit Festo},
pdfauthor={Steffen Sieder},
pdfsubject={Modellbasierte Funktionsentwicklung von},
pdfkeywords={},
pdftex,
bookmarks,
bookmarksopen,
breaklinks,
colorlinks=TRUE,
hyperindex=TRUE,
plainpages=FALSE,
pdfstartview=FitH,
bookmarksopenlevel=1,
bookmarksnumbered=true,
citecolor=black,
urlcolor=black,
linkcolor=black
]{hyperref}

% --------------------------
% autoref Namen (deutsch)
% --------------------------
\addto\extrasngerman{%
	\renewcommand{\sectionautorefname}{Abschnitt}%
	\renewcommand{\subsectionautorefname}{Abschnitt}%
	\renewcommand{\subsubsectionautorefname}{Abschnitt}%
}

% --------------------------
% Gleichungsnummerierung in Klammern (aus Vorlage)
% --------------------------
\makeatletter
\let\oldtheequation\theequation
\renewcommand\tagform@[1]{\maketag@@@{\ignorespaces#1\unskip\@@italiccorr}} %damit bei Benennung der Gleichung keine Doppelklammer entsteht
\renewcommand\theequation{(\oldtheequation)}
\makeatother

% --------------------------
% TikZ / Styles (auskommentiert in Vorlage) - hier nur als Kommentar
% --------------------------
% %\usepackage{tikz}
% %\usetikzlibrary{shapes,arrows,fit,calc,positioning,automata}
% %\tikzstyle{block} = [draw, fill=blue!20, rectangle, minimum height=3em, minimum width=\blockwidth]

% --------------------------
% Listings-Einstellungen (aus alter Preamble)
% --------------------------
\definecolor{dkgreen}{rgb}{0,0.6,0}
\definecolor{gray}{rgb}{0.5,0.5,0.5}
\definecolor{mauve}{rgb}{0.58,0,0.82}
\lstset{frame=none,
	aboveskip=3mm,
	belowskip=5mm,
	captionpos=b,
	showstringspaces=false,
	columns=flexible,
	basicstyle={\small\ttfamily},
	numbers=left,
	numberstyle=\tiny\color{gray},
	numbersep=5pt,
	keywordstyle=\color{blue},
	commentstyle=\color{dkgreen},
	stringstyle=\color{mauve},
	breaklines=true,
	breakatwhitespace=true,
	tabsize=3,
	inputpath = ./Skripte
}
\lstset{literate=%
	{Ö}{{\"O}}1
	{Ä}{{\"A}}1
	{Ü}{{\"U}}1
	{ß}{{\ss}}1
	{ü}{{\"u}}1
	{ä}{{\"a}}1
	{ö}{{\"o}}1
}
\renewcommand{\lstlistingname}{Skript}
\renewcommand{\lstlistlistingname}{Skriptverzeichnis}

% --------------------------
% Aufzählungs- und Absatz-Einstellungen sowie Inhalts-/Nummerierungstiefe
% --------------------------
\renewcommand{\labelenumi}{\arabic{enumi}.}

\setlength{\parindent}{0pt} % Einrücken nach Bildern verhindern
\setlength{\parskip}{12pt}	% globale Länge festlegen
\setcounter{tocdepth}{2}	% Inhaltsverzeichnis-Tiefe
\setcounter{secnumdepth}{2} % Nummerierungstiefe

% --------------------------
% Ende header.tex (vorläufig)
% --------------------------


% --------------------------
% Projektpakete (extern, bleiben als \usepackage)
% --------------------------
% Verzeichnisse, Titelblatt, util bleiben als separate Dateien/Pakete
\usepackage{Verzeichnisse}
\usepackage{Titelblatt}
\usepackage{util}

% --------------------------
% Literaturverzeichnis (biblatex)
% --------------------------
% https://www.overleaf.com/learn/latex/Biblatex_citation_styles
\usepackage[style=numeric, 
sorting=none,        % none --> nach reihenfolge der zitate im text
backend=bibtex,      % or backend=biber
bibencoding=utf8,    % .bib file encoding
maxbibnames=99       % give full list of authors (up to 99)
]{biblatex}
\setlength{\bibitemsep}{0.5\baselineskip} % Abstand zwischen zwei Quellen im Litver.
\bibliography{BibTeX_database_FSDR} % oder \addbibresource{Referenzen.bib} bei backend=biber

% (Optional: Beispiele für xpatch-Einsätze aus der alten Vorlage)
% \usepackage{xpatch}
% \xpretobibmacro{url}{\begingroup\color{red}}{}{}
% \xapptobibmacro{url}{\endgroup}{}

% --------------------------
% Seitenlayout / Kopf-/Fußzeile Einstellungen
% (Hinweis: viele dieser Grundfunktionen liefert classTimo bereits;
%  hier werden nur die projekt-spezifischen Anpassungen aus der alten Preamble übernommen)
% --------------------------

% Abstand vor/nach chapter-Überschrift
\renewcommand*\chapterheadstartvskip{\vspace{0cm}} % Abstand vor chapter Überschrift
\renewcommand*\chapterheadendvskip{\vspace{0cm}}   % Abstand nach chapter Überschrift

% Kopf-/Fußzeilen-Setup (scrlayer-scrpage verwendet classTimo bereits; wir überschreiben hier projekt-spezifisch)
\clearpairofpagestyles
\automark{chapter} % was soll in headmark stehen
\ihead[]{\headmark}  % i = links, c = mittig, o = rechts / headmark: Überschiften
\cfoot*{\pagemark} % [] Überschiftseiten, {} normale Seiten, * für beide
%\ohead*{\includegraphics[width=0.2\textwidth]{./Bilder/Stihl_header}}
\addtokomafont{pagehead}{\bfseries} % upshape = Kursiv, bfseries = Fett
\addtokomafont{pagenumber}{\bfseries}
% Auskommentierte zusätzliche Kopf-/Fuß-Optionen aus Vorlage (nur bei Bedarf aktivieren):
% %\pagestyle{scrheadings}
% %\clearscrheadfoot
% %\setheadsepline{0.5pt}
% %\setfootsepline{0.5pt}
% %\renewcommand*\chapterpagestyle{scrheadings}

% --------------------------
% (Auskommentierte) Verzeichnis-Anpassungen aus Vorlage
% Diese Einstellungen betreffen tocloft/tocbasic; da sie im Original auskommentiert waren,
% belasse ich sie hier kommentiert. Aktiviere nur falls benötigt.
% --------------------------
% %\renewcommand{\cftfigpresnum}{Abb. }	% Anpassen des Abbildungsverzeichnises
% %\renewcommand{\cfttabpresnum}{Tab. }
% %\renewcommand{\cftfigaftersnum}{:}
% %\renewcommand{\cfttabaftersnum}{:}
% %\setlength{\cftfignumwidth}{2cm}
% %\setlength{\cfttabnumwidth}{2cm}
% %\setlength{\cftfigindent}{0cm}
% %\setlength{\cfttabindent}{0cm}

% --------------------------
% Kurze Makros und Farben (aus alter Preamble)
% --------------------------
\newcommand\LM{\ensuremath{\mathit{LM}}}
\newcommand\IS{\ensuremath{\mathit{IS}}}

\usepackage{color}
\definecolor{FE_blue}{rgb}{0,0.5686,0.8627} % #0091DC
\definecolor{FE_blue2}{rgb}{0.8627,0.9216,0.9647} % #DCEBF6
\definecolor{FE_gray}{rgb}{0.8471,0.8627,0.8824} % 
\definecolor{AleeRed}{rgb}{0.5,0,0}

% --------------------------
% Mathematische Theoreme (aus Vorlage)
% --------------------------
% auskommentierte Varianten belassen; aktive Definition übernimmt shadethm-Umgebung
% %\usepackage{thm-autoref}
% %\newtheorem{satz}{Satz}
\newshadetheorem{satz}{Satz}
\newcommand{\shadesatzname}{Satz}
\newenvironment{thm}[1][]{%
	\definecolor{shadethmcolor}{FE_gray}%
	\definecolor{shaderulecolor}{rgb}{1.0,0.0,0.0}%
	\setlength{\shadeboxrule}{1pt}%
	\begin{sthm}[#1]%
	}{\end{sthm}}

% --------------------------
% hyperref (aus alter Vorlage) - am Ende der Pakete laden
% --------------------------
\usepackage[
pdftitle={Masterarbeit Festo},
pdfauthor={Steffen Sieder},
pdfsubject={Modellbasierte Funktionsentwicklung von},
pdfkeywords={},
pdftex,
bookmarks,
bookmarksopen,
breaklinks,
colorlinks=TRUE,
hyperindex=TRUE,
plainpages=FALSE,
pdfstartview=FitH,
bookmarksopenlevel=1,
bookmarksnumbered=true,
citecolor=black,
urlcolor=black,
linkcolor=black
]{hyperref}

% --------------------------
% autoref Namen (deutsch)
% --------------------------
\addto\extrasngerman{%
	\renewcommand{\sectionautorefname}{Abschnitt}%
	\renewcommand{\subsectionautorefname}{Abschnitt}%
	\renewcommand{\subsubsectionautorefname}{Abschnitt}%
}

% --------------------------
% Gleichungsnummerierung in Klammern (aus Vorlage)
% --------------------------
\makeatletter
\let\oldtheequation\theequation
\renewcommand\tagform@[1]{\maketag@@@{\ignorespaces#1\unskip\@@italiccorr}} %damit bei Benennung der Gleichung keine Doppelklammer entsteht
\renewcommand\theequation{(\oldtheequation)}
\makeatother

% --------------------------
% TikZ / Styles (auskommentiert in Vorlage) - hier nur als Kommentar
% --------------------------
% %\usepackage{tikz}
% %\usetikzlibrary{shapes,arrows,fit,calc,positioning,automata}
% %\tikzstyle{block} = [draw, fill=blue!20, rectangle, minimum height=3em, minimum width=\blockwidth]

% --------------------------
% Listings-Einstellungen (aus alter Preamble)
% --------------------------
\definecolor{dkgreen}{rgb}{0,0.6,0}
\definecolor{gray}{rgb}{0.5,0.5,0.5}
\definecolor{mauve}{rgb}{0.58,0,0.82}
\lstset{frame=none,
	aboveskip=3mm,
	belowskip=5mm,
	captionpos=b,
	showstringspaces=false,
	columns=flexible,
	basicstyle={\small\ttfamily},
	numbers=left,
	numberstyle=\tiny\color{gray},
	numbersep=5pt,
	keywordstyle=\color{blue},
	commentstyle=\color{dkgreen},
	stringstyle=\color{mauve},
	breaklines=true,
	breakatwhitespace=true,
	tabsize=3,
	inputpath = ./Skripte
}
\lstset{literate=%
	{Ö}{{\"O}}1
	{Ä}{{\"A}}1
	{Ü}{{\"U}}1
	{ß}{{\ss}}1
	{ü}{{\"u}}1
	{ä}{{\"a}}1
	{ö}{{\"o}}1
}
\renewcommand{\lstlistingname}{Skript}
\renewcommand{\lstlistlistingname}{Skriptverzeichnis}

% --------------------------
% Aufzählungs- und Absatz-Einstellungen sowie Inhalts-/Nummerierungstiefe
% --------------------------
\renewcommand{\labelenumi}{\arabic{enumi}.}

\setlength{\parindent}{0pt} % Einrücken nach Bildern verhindern
\setlength{\parskip}{12pt}	% globale Länge festlegen
\setcounter{tocdepth}{2}	% Inhaltsverzeichnis-Tiefe
\setcounter{secnumdepth}{2} % Nummerierungstiefe

% --------------------------
% Ende header.tex (vorläufig)
% --------------------------
